\textbf{Wedge.} (metal, NV, 2014; no injury)\newline \textit{``A wedge (block of rocks) came out of the wall along side of the haul road and struck Loader 377 while it was being used for high wall cleanup''} \\
\midrule
\textbf{Toppling.} (metal, NV, 2020; no injury)\newline \textit{``Multiple Geologic faults along the highwall contributed to a parallel joint orientation that dips steeply to the southwest. Bounding the instability to the north \& south vertical faults. These faults acted as ``release'' planes that resulted in a toppling deformation mechanism and movement of the highwall within this narrow sector.''} \\
\midrule
\textbf{Rotational slide.} (metal, MT, 2005; no injury)\newline \textit{``Due to high rain fall the Highwall was saturated with moisture and was being monitered.  Highwall activity indicated to evacuate employees from the area. Following our safety procedures employees were evacuated. No employee was injured or eqipment was damaged. There was no rock burst just a large slump of wet material from the highwall onto 5175 bench.''} \\
\midrule
\textbf{Collapse, mechanism unstated.} (stone, OR, 2020; days away and restricted work)\newline \textit{``Excavator was producing material from pit floor in front of lower bench of highwall. For unknown reason highwall gave way \& material made contact w/ excavator \& cab. Material entered cab \& EE received cut above lt eye at hairline on scalp. EE was taken to ER for check up \& 1st Aid on Cut. On 9/21/20 EE went back for x-rays which showed 6 cracked ribs \& a cracked sternum.''} \\
\midrule
\textbf{Rockfall.} (metal, NV, 2000; restricted work)\newline \textit{``WHILE BEING LOADED AT THE SHOVEL A LARGE ROCK ROLLED FROM THE HIGHWALL STRIKING THE REAR OF THE EMPLOYEE'S HAUL TRUCK. THE TRUCK ROCKED RESULTING IN MINOR NECK STRAIN. EMPLOYEE GIVEN 3 DAYS LIGHT DUTY BY PHYSICIAN.''} \\
